\chapter{General Introduction}
\label{ch:introduction}

The field of cosmology has advanced dramatically over the last few decades, driven by increasingly precise observations and crucial theoretical advancements, allowing us to construct a remarkably coherent picture of our cosmic history. Today, our understanding is encapsulated in the Standard Model of Cosmology, also known as the $\Lambda$CDM model, which describes the evolution of all cosmic components --- photons, neutrinos, atomic nuclei, electrons, and more --- from the earliest moments to the present, according to: (i) Einstein's theory of general relativity, which provides the geometric framework governing the dynamics of spacetime and its response to matter and energy; (ii) the Standard Model of particle physics, which describes the microphysical interactions of all known particles; and (iii) three new ingredients whose gravitational effects are well established but whose fundamental nature remains unknown.

The first new ingredient is \textit{inflation}, which posits an early period of nearly exponential expansion of the universe driven by a new hypothetical particle, the inflaton. This accelerated expansion ensures that when the process ends, the universe is sufficiently flat and uniform at the largest observable scales --- the very initial conditions that the subsequent hot Big Bang evolution requires. In addition, inflation provides a mechanism for generating the small density fluctuations that later evolve into the large structures observed today in the cosmic web. The second new ingredient is \textit{dark matter} (DM), a nearly collisionless and non-luminous form of matter, postulated to account for a multitude of gravitational effects observed across a wide range of scales --- from the formation and growth of large-scale structure to the dynamics of individual galaxies and clusters. Within the $\Lambda$CDM model, DM  accounts for about $80\%$ of the matter content in the universe, making it five times more prevalent than ordinary matter, the kind that makes up stars, planets, and all visible objects like us. Finally, the third new ingredient is the \textit{cosmological constant} $\Lambda$, a fixed energy density of the vacuum with negative pressure that drives the current accelerated expansion of the universe, first inferred from observations of distant supernovae in the late 1990s. While dark matter and ordinary matter together govern the formation and growth of structure, $\Lambda$ comes to dominate the total energy budget only at late times, accounting today for roughly 70\% of it. More generally, any component playing this role is referred to as \textit{dark energy}, and whether it is truly a cosmological constant or something that evolves dynamically over cosmic time is one of the central open questions in cosmology.


Within this simple yet powerful framework, after inflation the universe expanded and cooled from a hot, dense plasma of particles and radiation --- the hot Big Bang --- up to the point where the first light nuclei formed, just minutes after the Big Bang once the temperature dropped below a few MeV, in a process known as Big Bang Nucleosynthesis (BBN). The primary products of BBN are hydrogen, helium, and deuterium, along with trace amounts of heavier nuclei, and their predicted abundances depend sensitively on the number of relativistic species and the expansion rate at that epoch, making BBN one of the earliest indirect probes of the energy content of the universe. As cooling continued, roughly 380,000 years after the Big Bang, once the temperature dropped below about $0.3$ eV, electrons and protons were finally able to combine into neutral hydrogen, allowing photons --- which until then had been continuously scattering off the charged plasma --- to decouple and stream freely through space. These photons reach us today as a Cosmic Microwave Background (CMB) of radiation, carrying in its temperature and polarization patterns an imprint of the state of the universe at the moment of decoupling and of the physical processes that shaped it up to that point. Subsequently, the small density fluctuations seeded by inflation grew under gravity, with the first galaxies assembling roughly a billion years later and eventually giving rise to the large-scale distribution of galaxies and clusters we observe today, 13.8 billion years after the Big Bang.


The bedrock of this cosmological knowledge comes from observations of the CMB, which also constitutes the central observational focus of this thesis. In fact, its temperature and polarization anisotropies encode a wealth of information about the geometry and composition of the universe, the primordial conditions set during inflation, and the properties of all particle species present in the early universe. Further, secondary anisotropies --- imprinted on the CMB photons as they travel through the intervening matter and structure --- carry direct information about how those structures formed and evolved over cosmic time.